% =============================================================================
% 2.4 Positioning Statement — The Gap This Thesis Fills
% =============================================================================

\section{Positioning Statement --- The Gap This Thesis Fills}
\label{sec:positioning}

Each pillar has been studied independently. Resource contention theory is well-established \citep{nanda1991, arora2023}. Kubernetes overhead is measured at the infrastructure level \citep{casalicchio2019, liu2024}. Autoscaling behaviour is documented \citep{nguyen2020, tran2022, choi2021}. Fault isolation principles are understood architecturally \citep{cilic2023}.

What is missing---and what this thesis provides---is a study that:

\begin{enumerate}
  \item Connects resource contention and Kubernetes overhead at the \textbf{application-level workflow execution layer}, not the infrastructure level.
  \item Measures the \textbf{same workload} on \textbf{both architectures} across a range of concurrency levels.
  \item Identifies a specific \textbf{crossover threshold} where the migration becomes net beneficial.
  \item Uses \textbf{real pipeline orchestration workloads} (Dagster), not HTTP benchmarks or CPU stress tests.
  \item Runs on \textbf{real infrastructure} under controlled conditions, not simulations.
\end{enumerate}

No existing study provides engineers with a data-backed answer to ``at what workload level does the Kubernetes migration become worth it for pipeline orchestration systems?'' This thesis fills that gap.
